\section{Introduction}

Traditionally IT has had to choose less efficient means of storage for
certain classes of data (e.g. long-term data storage) for economic
reasons. To this end, legacy systems have used means such as tape or
optical to offload data from primary HDD storage for long-term archival.
However, such a trade-off becomes something of a faustian bargain, where
reduction in purchase price is traded for decreased efficiency
attributing to the manual nature of the medium. Today the dramatic
reduction in HDD cost \cite{ref1} coupled with MTTF rates greater than 100
years \cite{ref2} obviates the need for IT to make do with lesser means of
data storage. As these twin dynamics play out over time, HDD will
increasingly become the medium of choice for all data storage; removing
the need for IT to artificially create separate storage policies for
various types of data.

While reduced cost and increased reliability solves the problem of
physical medium, significant work remains to improve the lot of both IT
Operations (Ops) and the application developer. Long-term storage of
data, where time is measured in years to decades, introduces challenges
beyond those typically associated with short-term data storage.
Specifically: how can applications efficiently and effectively manage
the complexity of 1000's of millions (10\textsuperscript{9}) to millions
of millions (10\textsuperscript{12}) of files over extended periods of
time.

The continued increase in data creation is well documented. In
particular, there is substantial evidence that storage requirements for
one category of data is increasing at an increasing rate \cite{ref5,ref6}.
This class of data has several monikers, the more prevalent include:
fixed content, reference data, and unstructured data. Whatever the name,
this class of data is characterized as data that changes either
infrequently or not at all; e.g. medical imaging, archived documents and
email, etc.

To move forward in these areas requires advancements in software
methodology to complement the advances in hardware. This paper looks at
means of increasing abstraction from the perspective of the application
and IT Ops, as well as presenting a new model for addressing storage
based on the content of the data itself.
